%! AP Statistics — Unit 2, Lesson 2.8 — Group Activity
\documentclass[10pt]{article}
\usepackage{apstats-article}
\usepackage{apstats-boxes}

\begin{document}

\pageheader{Unit 2, Lesson 2.8}{Group Activity}
\namepartnerperiod

\vspace{0.12in}

\begin{tcolorbox}[colback=sky, colframe=navy, boxrule=0.9pt, arc=2mm,
    left=4mm, right=4mm, top=2.5mm, bottom=2.5mm]
\small\textit{A school nutritionist collects data on 40 students. For each student she records
daily screen time in hours ($x$) and a composite wellness score ($y$, 0--100 scale, higher is
better). Summary statistics: $\bar{x} = 4.2$ hr, $\bar{y} = 71$, $s_x = 1.8$ hr, $s_y = 12$,
$r = -0.75$.}
\end{tcolorbox}

\vspace{0.14in}

% ── TIER R ───────────────────────────────────────────────────────────────────
\begin{tcolorbox}[colback=redbg, colframe=black!40, boxrule=0.7pt, arc=2mm,
    left=4mm, right=4mm, top=2.5mm, bottom=2.5mm,
    title={\textbf{Tier R --- Remediate}}, fonttitle=\bfseries]
\small

\textbf{Use the formula reference below for each step.}

\vspace{4pt}
\begin{tabularx}{\linewidth}{@{} >{\bfseries}p{4cm} X @{}}
Slope formula & $b = r \cdot \dfrac{s_y}{s_x}$ \\[4pt]
Intercept formula & $a = \bar{y} - b\bar{x}$ \\
\end{tabularx}

\vspace{8pt}
\begin{enumerate}[leftmargin=1.6em, itemsep=0.16in]
  \item Substitute the values into the slope formula and compute $b$:\\[6pt]
        $b = \blank{2cm} \cdot \dfrac{\blank{2cm}}{\blank{2cm}} = \blank{2.5cm}$

  \item Substitute $b$, $\bar{x}$, and $\bar{y}$ into the intercept formula and compute $a$:\\[6pt]
        $a = \blank{2.5cm} - (\blank{2.5cm})(\blank{2.5cm}) = \blank{2.5cm}$

  \item Write the regression equation using the variable names (not $x$ and $y$):\\[6pt]
        \blank{10cm}

  \item Compute $r^2 = \blank{3cm}$. Circle the correct interpretation:

        \vspace{6pt}
        \textbf{(A)} About 56\% of students have a wellness score that is linear.\\
        \textbf{(B)} About 56\% of the variation in wellness score is explained by the regression on screen time.\\
        \textbf{(C)} The correlation between screen time and wellness is 56\%.
\end{enumerate}
\end{tcolorbox}

\vspace{0.12in}

% ── TIER A ───────────────────────────────────────────────────────────────────
\begin{tcolorbox}[colback=goldbg, colframe=black!40, boxrule=0.7pt, arc=2mm,
    left=4mm, right=4mm, top=2.5mm, bottom=2.5mm,
    title={\textbf{Tier A --- Approaching Proficiency}}, fonttitle=\bfseries]
\small

\begin{enumerate}[leftmargin=1.6em, itemsep=0.16in]
  \item Compute $b$ and $a$ from the summary statistics. Show all work.\\[8pt]
        \writeline\\[3pt]\writeline

  \item Write the regression equation with variable names. Then interpret the slope in context.\\[8pt]
        \writeline\\[3pt]\writeline\\[3pt]\writeline

  \item Compute $r^2$ and write the full AP-template interpretation.\\[8pt]
        \writeline\\[3pt]\writeline\\[3pt]\writeline

  \item Predict the wellness score for a student with 5 hours of daily screen time. Classify
        as interpolation or extrapolation. (The data ranged from 1 to 9 hours.)\\[8pt]
        \writeline\\[3pt]\writeline
\end{enumerate}
\end{tcolorbox}

\vspace{0.12in}

% ── TIER E ───────────────────────────────────────────────────────────────────
\begin{tcolorbox}[colback=greenbg, colframe=black!40, boxrule=0.7pt, arc=2mm,
    left=4mm, right=4mm, top=2.5mm, bottom=2.5mm,
    title={\textbf{Tier E --- Extension}}, fonttitle=\bfseries]
\small

\begin{enumerate}[leftmargin=1.6em, itemsep=0.16in]
  \item Complete all of Tier A above first.

  \item A new data point is added: a student with 14 hours of daily screen time and a wellness
        score of 85. Is this point likely to be \textit{influential}? Explain using the concepts
        of leverage and trend.\\[8pt]
        \writeline\\[3pt]\writeline\\[3pt]\writeline

  \item Without recomputing, predict whether adding this point would \textit{increase} or
        \textit{decrease} the slope, the intercept, and $r^2$. Justify each.\\[8pt]
        Slope: \writeline\\[3pt]
        Intercept: \writeline\\[3pt]
        $r^2$: \writeline\\[3pt]\writeline

  \item The nutritionist removes the point and recomputes: $r^2$ changes from 0.5625 to 0.48.
        Does this confirm or contradict your prediction? What does the change tell us about the
        point's influence?\\[8pt]
        \writeline\\[3pt]\writeline\\[3pt]\writeline
\end{enumerate}
\end{tcolorbox}

\end{document}
